% Book 3, Year 1, chapter 18 -- heaps and priority queues.
% Provenance (never printed): sources/book3/BRIEFS.md, chapter 18.
\chapter{Heaps and Priority Queues}\label{ch:b1:heaps}

In the waiting room of a hospital's emergency department, the next patient to be seen is not the one who came first,
but the one in the most urgent state. Patients arrive all the time, and a patient's state can worsen while waiting. A
queue serves in order of arrival; this room needs a structure that serves in order of priority, accepts new entries,
and changes a priority, each in a few steps however many people wait. The \omterm{def:b1:heaps:binary-heap}{binary heap} does this inside a plain \omterm{def:b1:pointers-and-arrays:array}{array},
with no \omterm{def:b1:pointers-and-arrays:pointer}{pointer} at all, and it also sorts in place and merges sorted lists.

\begin{recall}
The depth of a node is its distance from the root, the height of a tree the largest depth
(\cref{def:b1:induction-and-termination:height}). A
\omterm{def:b1:dynamic-arrays:dynamic-array}{dynamic array} grows by doubling, at amortised constant cost (\cref{ch:b1:dynamic-arrays}). A queue serves its elements
in order of arrival (\cref{ch:b1:linked-structures}).
\end{recall}

\section{The priority-queue type}\label{sec:b1:heaps:pq}

\begin{definition}[Priority queue]\label{def:b1:heaps:pq}
A \emph{priority queue}\index{priority queue} is an abstract type holding elements with keys, called priorities, with
the operations insert, find-min (the element of smallest key), extract-min (find it and remove it), and often
\omterm{def:b1:heaps:decrease-key}{decrease-key}.
\end{definition}

A max-priority queue serves the largest key first; exchanging $<$ and $>$ turns one into the other. Two simple
implementations put the cost on one side. An unsorted \omterm{def:b1:pointers-and-arrays:array}{array} inserts at its end in $O(1)$, but finding the minimum scans
all the elements. An \omterm{def:b1:pointers-and-arrays:array}{array} kept sorted, in decreasing order, extracts its last element in $O(1)$, but an insertion
shifts elements to make room. For $n$ insertions followed by $n$ extractions, both cost $\Theta(n^2)$:

\omoutput{code/bachelor-1/18-heaps/outputs/pq_naive.txt}{Comparisons for $n$ random insertions then $n$ extractions:
quadratic for the \omterm{def:b1:pointers-and-arrays:array}{arrays}, $n \log n$ for the heap of this chapter.}\label{out:b1:heaps:naive}

\section{The binary heap in an array}\label{sec:b1:heaps:array}

\begin{definition}[Complete binary tree]\label{def:b1:heaps:complete}
A \omterm{def:b1:induction-and-termination:binary-tree}{binary tree} of height $h$ is \emph{complete}\index{complete binary tree} when its depths $0, \dots, h - 1$ are full
(depth $d$ holds $2^d$ nodes) and the nodes at depth $h$ are as far left as possible.
\end{definition}

Numbered level by level, from left to right, starting at 0, the nodes of a \omterm{def:b1:heaps:complete}{complete} tree fill the indices $0, 1,
\dots, n - 1$ with no gap: the tree can live in an \omterm{def:b1:pointers-and-arrays:array}{array}, and the links are computed instead of stored.

\begin{proposition}[Index arithmetic]\label{prop:b1:heaps:index}
In a \omterm{def:b1:heaps:complete}{complete binary tree} numbered from 0, the children of node $i$ are $2i + 1$ and $2i + 2$ (when they are below $n$),
and the parent of node $i > 0$ is $\lfloor (i - 1)/2 \rfloor$. A \omterm{def:b1:heaps:complete}{complete} tree of $n \geq 1$ nodes has height
$\lfloor \log_2 n \rfloor$.
\end{proposition}

\begin{proof}
Depth $d$ holds the indices $2^d - 1$ to $2^{d+1} - 2$. Let $i = 2^d - 1 + j$ be the $j$-th node of depth $d$ ($0 \leq j <
2^d$). Its children are the nodes $2j$ and $2j + 1$ of depth $d + 1$, with indices $2^{d+1} - 1 + 2j = 2i + 1$ and $2i +
2$. Conversely, both $2i + 1$ and $2i + 2$ give $\lfloor (k - 1)/2 \rfloor = i$. For the height, see
\cref{exo:b1:heaps:9}.
\end{proof}

\begin{definition}[Binary heap, heap property]\label{def:b1:heaps:binary-heap}
A \emph{binary heap}\index{binary heap} is a \omterm{def:b1:heaps:complete}{complete binary tree} stored in an \omterm{def:b1:pointers-and-arrays:array}{array}, with the \emph{heap
property}\index{heap property}: the key of every node is at most the keys of its children (a min-heap), or at least
them (a max-heap).
\end{definition}

By transitivity, the key of a node is at most every key of its \omterm{def:b1:induction-and-termination:rooted-tree}{subtree}, and the root holds the minimum: find-min costs
$O(1)$. A heap is much less ordered than a search tree: siblings are in no particular order, and the \omterm{def:b1:pointers-and-arrays:array}{array} is not
sorted. That is what makes it cheap to maintain.

\begin{omfigure}
\begin{tikzpicture}[level distance=9mm, n/.style={draw, circle, minimum size=6.5mm, font=\small, inner sep=1pt},
  level 1/.style={sibling distance=26mm}, level 2/.style={sibling distance=13mm}, >={Stealth[length=1.6mm]}]
\node[n] {1}
  child {node[n] {4} child {node[n] {7}} child {node[n] {8}}}
  child {node[n] {2} child {node[n] {9}} child {node[n] {3}}};
\node[font=\tiny] at (0.45,0.3) {0};
\begin{scope}[xshift=4.3cm, yshift=-1.2cm, c/.style={draw, minimum width=7mm, minimum height=6mm, font=\small}]
\foreach \v/\i in {1/0, 4/1, 2/2, 7/3, 8/4, 9/5, 3/6} {\node[c] (a\i) at (\i*0.7,0) {\v};
  \node[font=\tiny] at (\i*0.7,-0.5) {\i};}
\draw[->, omCode, thick] (a1.north) to[bend left=50] (a3.north);
\draw[->, omCode, thick] (a1.north) to[bend left=60] (a4.north);
\draw[->, omThm, thick] (4.2,-0.7) to[bend left=30] (1.4,-0.7);
\node[font=\scriptsize, text=omCode, anchor=west] at (3.6,0.9) {children of 1: 3 and 4};
\node[font=\scriptsize, text=omThm, anchor=west] at (4.6,-0.9) {parent of 6: $\lfloor 5/2 \rfloor = 2$};
\end{scope}
\end{tikzpicture}
\par\smallskip{\small The min-heap built by \cref{out:b1:heaps:demo}, as a tree and as the \omterm{def:b1:pointers-and-arrays:array}{array} that stores it. Every
key is at most its children's keys; the arrows are computed from the indices, not stored.}
\end{omfigure}

\section{Sifting up and down}\label{sec:b1:heaps:sift}

\begin{definition}[Sift up, sift down, heapify]\label{def:b1:heaps:sift}
To \emph{sift up}\index{sift up} a node is to exchange it with its parent as long as its key is smaller. To \emph{sift
down}\index{sift down} a node is to exchange it with its smaller child as long as that child's key is smaller. To
\emph{heapify}\index{heapify} an \omterm{def:b1:pointers-and-arrays:array}{array} is to make it a heap by sifting down every node that has a child, from the last
one back to the root.
\end{definition}

\omcode{code/bachelor-1/18-heaps/listings/heap.c}{13}{33}{Sifting up and down, with counters}\label{lst:b1:heaps:sift}

Each loop keeps an invariant. For sift-up: \emph{the \omterm{def:b1:pointers-and-arrays:array}{array} is a heap except that node $i$ may be smaller than its
parent}. If it is not, the \omterm{def:b1:pointers-and-arrays:array}{array} is a heap. If it is, the exchange puts the smaller key above both former children of the
parent, which were at least the old parent's key, and only the parent's position may now be wrong. For sift-down:
\emph{the \omterm{def:b1:pointers-and-arrays:array}{array} is a heap except that node $i$ may be larger than one of its children}; exchanging it with the smaller
child puts above the other child a key at most that child's key. Each step moves one level, so both cost $O(\log n)$.

Insertion places the new key as a \omterm{def:b1:induction-and-termination:rooted-tree}{leaf} at the end of the \omterm{def:b1:pointers-and-arrays:array}{array}, which keeps the tree \omterm{def:b1:heaps:complete}{complete}, and sifts it up.
Extraction saves the root, moves the last \omterm{def:b1:induction-and-termination:rooted-tree}{leaf} to the root, and sifts it down.

\omcodes{code/bachelor-1/18-heaps/listings/heap.c}{41-58}{Insertion and extraction}\label{lst:b1:heaps:push-pop}

\omoutput{code/bachelor-1/18-heaps/outputs/heap_demo.txt}{The \omterm{def:b1:pointers-and-arrays:array}{array} after each insertion and
extraction.}\label{out:b1:heaps:demo}

\begin{omfigure}
\begin{tikzpicture}[level distance=8mm, n/.style={draw, circle, minimum size=6mm, font=\small, inner sep=1pt},
  m/.style={n, fill=omThm!25}, level 1/.style={sibling distance=20mm}, level 2/.style={sibling distance=10mm},
  level 3/.style={sibling distance=8mm}, >={Stealth[length=1.6mm]}]
\node[m] {10}
  child {node[n] {3} child {node[n] {6} child {node[n] {8}} child {node[n] {9}}} child {node[n] {4}}}
  child {node[n] {2} child {node[n] {5}} child {node[n] {7}}};
\draw[->, thick] (2.1,-1.3) -- (2.9,-1.3);
\begin{scope}[xshift=5cm]
\node[n] {2}
  child {node[n] {3} child {node[n] {6} child {node[n] {8}} child {node[n] {9}}} child {node[n] {4}}}
  child {node[m] {10} child {node[n] {5}} child {node[n] {7}}};
\draw[->, thick] (2.1,-1.3) -- (2.9,-1.3);
\end{scope}
\begin{scope}[xshift=10cm]
\node[n] {2}
  child {node[n] {3} child {node[n] {6} child {node[n] {8}} child {node[n] {9}}} child {node[n] {4}}}
  child {node[n] {5} child {node[m] {10}} child {node[n] {7}}};
\end{scope}
\end{tikzpicture}
\par\smallskip{\small Extract-min on the heap $[1, 3, 2, 6, 4, 5, 7, 8, 9, 10]$: the last \omterm{def:b1:induction-and-termination:rooted-tree}{leaf}, 10, replaces the root 1
(left); it is exchanged with its smaller child 2 (middle), then with 5 (right), where both its \omterm{def:b1:induction-and-termination:rooted-tree}{subtrees} are empty. Two
exchanges, four comparisons.}
\end{omfigure}

\section{Building a heap in linear time}\label{sec:b1:heaps:build}

A heap of $n$ given keys can be built by $n$ insertions, at cost $O(n \log n)$; on decreasing keys, each new key climbs
to the root, and the cost is $\Theta(n \log n)$. \omterm{def:b1:heaps:sift}{Heapify} does better: it sifts down the nodes from the last parent,
$\lfloor n/2 \rfloor - 1$, back to the root; when node $i$ is sifted down, both its \omterm{def:b1:induction-and-termination:rooted-tree}{subtrees} are already heaps (they were
treated before it), so it alone may be out of place.

\begin{proposition}[Heapify is linear]\label{prop:b1:heaps:build}
\omterm{def:b1:heaps:sift}{Heapify} makes at most $\sum_i h_i < n$ exchanges, where $h_i$ is the height of node $i$.
\end{proposition}

\begin{proof}
Sifting node $i$ down makes at most $h_i$ exchanges. A node of height at least $h$ has a descendant $h$ levels below it,
on its leftmost path, which is node $(i + 1)2^h - 1$; so the nodes of height at least $h$ are those with $i + 1 \leq
n/2^h$, and there are $\lfloor n/2^h \rfloor$ of them. Each node of height $h_i$ is counted once for each $h = 1, \dots,
h_i$: $\sum_i h_i = \sum_{h \geq 1} \lfloor n/2^h \rfloor < n(\frac12 + \frac14 + \cdots) = n$.
\end{proof}

The problem shows that the sum is exactly $n - \nu(n)$, where $\nu(n)$ is the number of ones in the binary writing of
$n$. On decreasing keys every node goes down to a \omterm{def:b1:induction-and-termination:rooted-tree}{leaf}, and the count is reached exactly:

\omoutput{code/bachelor-1/18-heaps/outputs/build.txt}{Exchanges to build a heap of decreasing keys: bottom-up below $n$,
top-down growing like $n \log_2 n$.}\label{out:b1:heaps:build}

\begin{center}\small
\begin{tabular}{lcccc}
\hline
depth of the nodes (tree of 15 nodes) & 0 & 1 & 2 & 3 \\
\hline
number of nodes & 1 & 2 & 4 & 8 \\
longest sift-down (their height) & 3 & 2 & 1 & 0 \\
exchanges at most & 3 & 4 & 4 & 0 \\
\hline
\end{tabular}
\end{center}

\noindent Most nodes are near the bottom and move little; the few that may move far are near the top. In total $3 + 4 +
4 = 11 = 15 - \nu(15)$ exchanges at most. Insertion one by one has it the other way round: the many leaves may each climb
the whole height.

\section{Heapsort and other uses}\label{sec:b1:heaps:uses}

\begin{method}[Heapsort]\label{met:b1:heaps:heapsort}
\emph{Heapsort}\index{heapsort} sorts an \omterm{def:b1:pointers-and-arrays:array}{array} in increasing order in place: \omterm{def:b1:heaps:sift}{heapify} it as a max-heap; then, for $end =
n - 1$ down to 1, exchange the root (the maximum of $a[0..end]$) with $a[end]$, and sift the new root down in
$a[0..end)$. The invariant: $a[0..end]$ is a max-heap, and $a[end + 1..n)$ holds the largest keys, sorted, all at least
those of the heap.
\end{method}

\omcode{code/bachelor-1/18-heaps/listings/heapsort.c}{17}{24}{\omterm{met:b1:heaps:heapsort}{Heapsort}}\label{lst:b1:heaps:heapsort}

\begin{omfigure}
\begin{tikzpicture}[c/.style={draw, minimum width=7mm, minimum height=6mm, font=\small}]
\foreach \row/\vals/\k/\lab in {0/{9,5,8,4,2,3,7,1}/0/{max-heap},
  1/{8,5,7,4,2,3,1,9}/1/{1 extraction}, 2/{7,5,3,4,2,1,8,9}/2/{2 extractions},
  3/{5,4,3,1,2,7,8,9}/3/{3 extractions}} {
  \node[font=\small, anchor=east] at (-0.6,-\row*0.75) {\lab};
  \foreach \v [count=\j from 0] in \vals {
    \pgfmathtruncatemacro{\sorted}{\j >= 8 - \k ? 1 : 0}
    \ifnum\sorted=1 \node[c, fill=omDef!20] at (\j*0.7,-\row*0.75) {\v};
    \else \node[c] at (\j*0.7,-\row*0.75) {\v}; \fi}}
\node[font=\scriptsize, anchor=west] at (5.8,-2.25) {heap part, then the sorted part (shaded)};
\end{tikzpicture}
\par\smallskip{\small \omterm{met:b1:heaps:heapsort}{Heapsort} of $[5, 2, 8, 1, 9, 3, 7, 4]$: the \omterm{def:b1:pointers-and-arrays:array}{array} after \omterm{def:b1:heaps:sift}{heapify} and after each of the first three
extractions. The sorted part grows from the right; the heap part shrinks.}
\end{omfigure}

\omoutput{code/bachelor-1/18-heaps/outputs/heapsort.txt}{\omterm{met:b1:heaps:heapsort}{Heapsort} of random \omterm{def:b1:pointers-and-arrays:array}{arrays}: comparisons below $2 n \log_2
n$.}\label{out:b1:heaps:heapsort}

\omterm{met:b1:heaps:heapsort}{Heapsort} runs in $\Theta(n \log n)$ in the worst case and uses $O(1)$ extra memory, which no sort of the previous
chapters achieves together; merge sort, of \cref{ch:b1:sorting}, needs a second \omterm{def:b1:pointers-and-arrays:array}{array}. It is not
\emph{stable}: elements with equal keys can come out in another order than they came in, since the exchanges move
elements across long distances.

\omoutput{code/bachelor-1/18-heaps/outputs/stability.txt}{Heapsorting patients by urgency: Ben arrived before Finn, but
leaves after him.}\label{out:b1:heaps:stability}

A heap of size $k$ selects the $k$ largest of $n$ values in one pass, in $O(n \log k)$ time and $O(k)$ memory: its root
is the smallest of the $k$ kept, the one a new value must beat.

\omcode{code/bachelor-1/18-heaps/listings/topk.c}{6}{26}{The $k$ longest loans of the town library, with a heap of
size $k$}\label{lst:b1:heaps:topk}

\omoutput{code/bachelor-1/18-heaps/outputs/topk.txt}{The eight longest returned loans.}\label{out:b1:heaps:topk}

\begin{definition}[Decrease-key]\label{def:b1:heaps:decrease-key}
The operation \emph{decrease-key}\index{decrease-key} lowers the key of an element already in a \omterm{def:b1:heaps:pq}{priority queue}.
\end{definition}

In a heap, a smaller key can only move up, by a sift-up from the element's position, in $O(\log n)$. The difficulty is to
find that position: the heap is not a search structure. An \emph{indexed heap} stores the items $0, \dots, m - 1$ and an
\omterm{def:b1:pointers-and-arrays:array}{array} \code{pos} that every move updates, so that the position of an item is read in $O(1)$. Shortest-path algorithms,
in \cref{ch:b1:shortest-paths}, use it. The alternative, \emph{lazy deletion}, inserts the element again with its new
key and leaves the old entry in the heap, to be skipped when it reaches the root: simpler, at the price of a larger heap.

\omcodes{code/bachelor-1/18-heaps/listings/iheap.c}{10-17,44-47}{Sift-up in an indexed heap, and
\omterm{def:b1:heaps:decrease-key}{decrease-key}}\label{lst:b1:heaps:iheap}

\omoutput{code/bachelor-1/18-heaps/outputs/er_demo.txt}{An emergency room with \omterm{def:b1:heaps:decrease-key}{decrease-key}; urgency 1 is the
highest.}\label{out:b1:heaps:er}

\omimg{book3/ai/b1-heaps-emergency-room.jpg}{0.55\linewidth}{At the triage desk of an emergency department,
patients are seen by urgency, not by order of arrival, and an urgency can rise while a patient waits: a \omterm{def:b1:heaps:pq}{priority
queue} with \omterm{def:b1:heaps:decrease-key}{decrease-key}.}

\section{Exercises}\label{sec:b1:heaps:exercises}

\begin{exercise}[$\star$]\label{exo:b1:heaps:1}
Give the \omterm{def:b1:pointers-and-arrays:array}{array} of a min-heap after inserting 5, 3, 8, 1, 4, 2 in this order into an empty heap.
\end{exercise}

\begin{exercise}[$\star$]\label{exo:b1:heaps:2}
Give the \omterm{def:b1:pointers-and-arrays:array}{array} after one extract-min on the heap $[1, 3, 2, 5, 4, 8]$.
\end{exercise}

\begin{exercise}[$\star$]\label{exo:b1:heaps:3}
In a heap numbered from 0, give the parent, the children and the depth of node 13.
\end{exercise}

\begin{exercise}[$\star$]\label{exo:b1:heaps:4}
Which of the \omterm{def:b1:pointers-and-arrays:array}{arrays} $[2, 5, 3, 9, 6, 4, 8, 10, 7]$ and $[1, 4, 2, 6, 5, 3]$ are min-heaps?
\end{exercise}

\begin{exercise}[$\star\star$]\label{exo:b1:heaps:5}
This sift-down compares the node with its left child only. Explain the output, and give the smallest heap on which one
extraction already goes wrong.
\omcode{code/bachelor-1/18-heaps/listings/exo_sift_bug.c}{4}{9}{Sift-down with one comparison}
\omoutput[2-2]{code/bachelor-1/18-heaps/outputs/exo_sift_bug.txt}{Seven extractions from the heap $[1, 5, 2, 6, 7, 3,
4]$.}
\end{exercise}

\begin{exercise}[$\star\star$]\label{exo:b1:heaps:6}
These formulas come from a text that numbers the \omterm{def:b1:pointers-and-arrays:array}{array} from 1. Explain each wrong line of the output, and correct the
formulas.
\omcode{code/bachelor-1/18-heaps/listings/exo_index_bug.c}{4}{5}{Parent and left child}
\omoutput[2-6]{code/bachelor-1/18-heaps/outputs/exo_index_bug.txt}{On the heap $[1, 3, 2, 7, 4]$.}
\end{exercise}

\begin{exercise}[$\star\star$]\label{exo:b1:heaps:7}
Write a max-heap whose elements are of any type, given by their size in bytes and a comparison function as for
\code{qsort}, and test it on numbers and on structures.
\end{exercise}

\begin{exercise}[$\star\star$]\label{exo:b1:heaps:8}
Write \code{kth_smallest(xs, n, k)} in time $O(n \log k)$ with a heap of size $k$, and test it against sorting.
\end{exercise}

\begin{exercise}[$\star\star$]\label{exo:b1:heaps:9}
Prove that a \omterm{def:b1:heaps:complete}{complete binary tree} with $n \geq 1$ nodes has height $\lfloor \log_2 n \rfloor$.
\end{exercise}

\begin{exercise}[$\star\star\star$]\label{exo:b1:heaps:10}
Prove by induction on $k$ that in a perfect \omterm{def:b1:induction-and-termination:binary-tree}{binary tree} of height $k - 1$, with $n = 2^k - 1$ nodes, the heights of the
nodes add up to $n - k$.
\end{exercise}

\begin{exercise}[$\star\star\star$]\label{exo:b1:heaps:11}
Maintain the median of a stream of numbers with two heaps, a max-heap of the lower half and a min-heap of the upper
half, so that each new number costs $O(\log n)$. Test it against sorting.
\end{exercise}

\begin{exercise}[$\star\star\star$]\label{exo:b1:heaps:12}
Prove that heapsort makes at most $2n \log_2 n$ key comparisons for $n \geq 2$, and compare with
\cref{out:b1:heaps:heapsort}.
\end{exercise}

\section{Problem: Merging the branch libraries' lists}\label{sec:b1:heaps:problem}

\begin{problem}[{Weekend problem --- Merging the branch libraries' lists}]\label{pb:b1:heaps:1}
The town library has $k$ branches; each keeps the list of the loans returned to it, sorted by day of return, and the
town wants one sorted list of all $n = 378$ returns. A program assigns the returned loans to $k$ branches (by book number
modulo $k$), then merges the $k$ sorted lists in two ways: by comparing the heads of all the lists at each step, and
with a min-heap that holds the head of each list. It counts the comparisons and checks the merged list.
\omoutput{code/bachelor-1/18-heaps/outputs/kmerge.txt}{Merging $k$ sorted lists of return days.}\label{out:b1:heaps:kmerge}
\omoutput{code/bachelor-1/18-heaps/outputs/heights_sum.txt}{Sums of heights in heaps of $n$
nodes.}\label{out:b1:heaps:heights}

\needspace{4\baselineskip}
\medskip\noindent\textbf{Part I --- Merging by scanning.}
\begin{enumerate}
\item Show that the scanning merge makes at most $(k - 1)n$ comparisons.
\item Compare with the table for $k = 32$. Why is the count smaller than the bound?
\item Why do both methods make the same number of comparisons for $k = 2$?
\item Another method merges the lists two by two, then the results two by two, and so on. What does it cost, for $k$ a
power of 2?
\item Why does the merge need only one value from each list at a time, and why does that matter for long lists stored in
files?
\item Prove that the merged list is sorted, whatever method chooses the smallest head.
\end{enumerate}

\needspace{4\baselineskip}
\medskip\noindent\textbf{Part II --- A heap of the heads.}
\begin{enumerate}[resume]
\item State the invariant of the heap during the merge.
\item Prove that the root is the smallest value not yet output.
\item Why is replacing the root and sifting it down better than an extract-min followed by an insertion?
\item Show that each output costs at most $2 \lfloor \log_2 k \rfloor$ comparisons, and bound the total.
\item Compare the bound with the table for $k = 32$.
\item The program stores $\textit{day} \times 64 + \textit{branch}$ in the heap. Why, and what does it assume about $k$?
\end{enumerate}

\needspace{4\baselineskip}
\medskip\noindent\textbf{Part III --- Building the heap.}
\begin{enumerate}[resume]
\item The merge builds its heap of $k$ heads by $k$ insertions. Would \omterm{def:b1:heaps:sift}{heapify} change the total cost noticeably?
\item From \cref{out:b1:heaps:build}, give the exchanges per key of the top-down build on decreasing keys, and relate
them to $\log_2 n$.
\item Why do decreasing keys make every node of the bottom-up build sink to a \omterm{def:b1:induction-and-termination:rooted-tree}{leaf}?
\item Check the bottom-up column against $n - \nu(n)$ for $n = 1023$ and $n = 1\,048\,575$.
\item Why does each sift-down of \omterm{def:b1:heaps:sift}{heapify} make at most two comparisons per exchange, plus two? Bound the comparisons of
\omterm{def:b1:heaps:sift}{heapify}.
\item \omterm{met:b1:heaps:heapsort}{Heapsort} begins with \omterm{def:b1:heaps:sift}{heapify}. Which part of heapsort dominates its cost?
\end{enumerate}

\needspace{4\baselineskip}
\medskip\noindent\textbf{Part IV --- The sum of heights.}
\begin{enumerate}[resume]
\item Show that node $i$ of a heap of $n$ nodes has height $\lfloor \log_2 (n/(i + 1)) \rfloor$.
\item Deduce that the sum of the heights is $S(n) = \sum_{h \geq 1} \lfloor n/2^h \rfloor$.
\item Prove that $\sum_{h \geq 1} \lfloor 2^j/2^h \rfloor = 2^j - 1$, and that $\lfloor n/2^h \rfloor$ is the sum of these
quantities over the bits $2^j$ of $n$.
\item Deduce $S(n) = n - \nu(n)$, and check it on $n = 1000$ in \cref{out:b1:heaps:heights}.
\item How far is $n - \nu(n)$ from the bound $n$ of \cref{prop:b1:heaps:build}?
\item State the named result of the problem, \emph{the heap identity}, with what it gives for \omterm{def:b1:heaps:sift}{heapify}.
\end{enumerate}
\end{problem}
